\documentclass[10pt]{article}
\usepackage[utf8]{inputenc}
\usepackage[T1]{fontenc}
\usepackage[margin=0.75in,top=0.6in,bottom=0.6in]{geometry}
\usepackage{microtype}
\usepackage{amsmath,amssymb,amsfonts,mathtools}
\usepackage{booktabs,array}
\usepackage{multicol}
\usepackage{graphicx}
\usepackage{xcolor}
\usepackage{enumitem}
\usepackage{tcolorbox}
\usepackage{fancyhdr}
\usepackage[colorlinks=true,linkcolor=blue,citecolor=blue,urlcolor=blue,hidelinks]{hyperref}
\usepackage[nameinlink,capitalize]{cleveref}

\setlength{\parskip}{1pt}
\setlength{\parindent}{0pt}
\setlength{\columnsep}{14pt}

\pagestyle{fancy}
\fancyhf{}
\fancyhead[L]{\small\color{blue!60!black}Project summary for feedback}
\fancyhead[R]{\small\color{blue!60!black}Invariance vs robustness}
\fancyfoot[R]{\thepage}
\renewcommand{\headrulewidth}{0pt}
\renewcommand{\arraystretch}{0.95}

\usepackage{titlesec}
\titlespacing*{\section}{0pt}{4pt}{2pt}
\titleformat*{\section}{\normalsize\bfseries}

\newcommand{\etaL}{\ensuremath{\eta/L}}

\begin{document}
\thispagestyle{fancy}

{\small\color{blue!60!black}Technical University of Munich \hspace{0.3em}$\vert$\hspace{0.3em} Practical course, extended to a paper \hspace{0.3em}$\vert$\hspace{0.3em} Supervisor: M.\ Sabanayagam}

\vspace{0.3em}
{\Large\bfseries Invariance vs Robustness: What Orders a Model's Adversarial Robustness\par}
\vspace{0.1em}
{\normalsize A threat-matched margin-to-Lipschitz account, from small classifiers to frozen foundation-model encoders\par}
\vspace{0.3em}
{\normalsize\bfseries Anass Al Ammiri}\hfill {\small July 19, 2026}

\vspace{0.3em}
{\small Current status: a combined paper at roughly 23 pages with a central concept figure and five empirical regimes, plus a full proofs-and-details appendix; code and result files reproduce every number. A pitch deck and an internal reader script accompany this summary. This is a two-page summary for feedback, not a submission plan.}
\vspace{0.3em}

\begin{tcolorbox}[colback=blue!3,colframe=blue!40!black,boxrule=0.5pt,arc=2pt,left=4pt,right=4pt,top=3pt,bottom=3pt]
\small
\textbf{Summary.} Whether shift-invariance helps or hurts adversarial robustness is contested: Ge et al.\ (2021) prove a fully shift-invariant classifier can be far less robust than a fully connected one, while anti-aliasing and equivariance work (Zhang 2019; Saha and Gokhale 2024; Wang 2025) reports the opposite correlation. Every ``helps'' result, though, used a gradient attack (FGSM or PGD), never a parameter-free strong attack. We organise both sides around one quantity. A shift-invariant linear classifier separates classes only through their projection onto the invariant subspace, with an exact margin identity (the distance between the projected convex hulls, Proposition~1). The surviving robustness is summarised by a threat-matched, scale-free margin-to-Lipschitz ratio \etaL, a Lipschitz--margin certificate (Hein 2017; Tsuzuku 2018) that we complete to a \emph{two-sided bracket} $\eta/L\le r\le\eta/\alpha$ with condition number $\kappa=L/\alpha\in[1.6,2.4]$ (Proposition~2). Across capacity-matched dissections and, importantly, within single frozen CLIP encoders where training cannot confound, \etaL{} orders the adversarial robust radius while shift-consistency does not; on modern encoders shift-consistency is saturated ($0.963$--$0.988$) and cannot rank at all. The ``invariance helps'' correlation is a \emph{weak-attack artifact}: it appears under FGSM and vanishes under AutoAttack. To rank already-robust models, where \etaL{} no longer separates, the \emph{shape} of the input gradient (an $\ell_1/\ell_2$ anisotropy) tracks robustness across thirty RobustBench models (Spearman $-0.79$), through a proposed spread-over-dimensions deficit.
\end{tcolorbox}

\vspace{0.2em}
\begin{multicols}{2}

\section*{1\quad Problem and motivation}
Image classifiers are often called shift-invariant when circular translations of the input do not change the prediction. The published evidence on what this does to adversarial robustness is mixed. Ge et al.\ (2021) show a fully shift-invariant network can be fooled by an $\ell_2$ perturbation of order $1/\sqrt{d}$, where a fully connected one needs order $1$. Anti-aliased pooling and equivariance work (Zhang 2019; Saha and Gokhale 2024; Grabinski 2023; Wang 2025) reports that operators which increase shift-invariance also increase robustness. A separate line proves a trade-off between geometric and adversarial robustness (Kamath et al.\ 2021; Tram\`er et al.\ 2020). The practitioner is left without a rule for when invariance is safe.

\textbf{Two under-tested assumptions.} The ``helps'' line only ever measured robustness with gradient attacks (FGSM/PGD), which over-report robustness, and it never tested whether the shift-consistency \emph{metric} actually orders robustness across models. Both gaps are testable.

\textbf{The gap.} There was no account that (i) says what discriminative signal survives a given invariance, (ii) turns that into a robustness certificate with a converse, and (iii) is validated against a standardized strong attack at matched capacity and, ideally, inside a single frozen model. This project supplies all three.

\section*{2\quad The quantity and its theory}
\textbf{Proposition 1 (invariant margin).} For a finite group acting orthogonally, an invariant linear classifier sees the data only through the group-average projector $P_G$; its maximum $\ell_2$ margin is $\tfrac12\,\mathrm{dist}(\mathrm{conv}\,P_GX_+,\mathrm{conv}\,P_GX_-)$, the separation of the class averages. The cyclic-shift case is the DC projection and recovers Ge's $2/\sqrt d$ collapse, now for any finite separable dataset, not one image.

\textbf{The criterion.} The surviving robustness is summarised by $\etaL$: the mean margin $\eta$ over the mean dual-norm gradient $L$ on clean, correctly classified images. We \emph{threat-match} the norm to the attack ($q{=}1$ for $\ell_\infty$, $q{=}2$ for $\ell_2$; the mismatched norm gives the wrong sign), and note it is \emph{scale-free}, so robustness attaches to the ratio, not to margin or Lipschitz alone. No attack is needed to compute it.

\textbf{Proposition 2 (two-sided bracket).} With an $L$-Lipschitz and $\alpha$-co-Lipschitz condition toward the boundary, $\eta/L\le r\le\eta/\alpha$, so the true radius is trapped, not merely lower-bounded; the gap is $\kappa=L/\alpha$, exact ($\kappa{=}1$) for linear invariant features and $\kappa\in[1.6,2.4]$ empirically. This is the converse the one-sided Lipschitz--margin certificate lacked.

\textbf{Further theory (appendix).} A convolution--square--pool layer computes power-spectrum coordinates; a degree-three bispectrum invariant separates the phase-coded classes the power spectrum cannot. On orbit-closed data invariance is \emph{margin-free} (invariant and unconstrained max-margin values coincide); it helps only on near-orthogonal cluster geometry, where a trajectory-level selection theorem is the open lazy-regime question.

\section*{3\quad Main empirical evidence}
\textbf{\etaL{} orders robustness; shift-consistency does not.} Across settings, the threat-matched ratio tracks the robust radius or AutoAttack robust accuracy, while shift-consistency is uncorrelated or anti-correlated.

\vspace{2pt}
\begin{tcolorbox}[colback=gray!4,colframe=gray!50,boxrule=0.4pt,arc=2pt,left=3pt,right=3pt,top=2pt,bottom=2pt]
\footnotesize
\begin{tabular}{@{}l@{\ \ }r@{\ \ }r@{}}
\toprule
setting (attack / metric) & \etaL & consist. \\
\midrule
CIFAR-10 standard (min-norm $r$)   & $+0.998$ & $-0.33$\,n.s. \\
PreActResNet-18 AT (AA $\ell_\infty$) & $+0.91$ & $-0.81$ \\
ImageNet-100 AT (AA $\ell_\infty$)    & $+0.90$ & $-0.87$ \\
16 CLIP encoders, per image ($r$)  & $+0.74$ to $+0.96$ & $\approx 0$ \\
\bottomrule
\end{tabular}\\[2pt]
Pearson or Spearman as appropriate. The per-image rows are \emph{within} a single frozen encoder, so training is not the confound.
\end{tcolorbox}
\vspace{1pt}

\textbf{The ``helps'' effect is a weak-attack artifact.} Re-running the ``helps'' regime (standard training, FGSM/PGD) on our operators and on Saha (2024) and Wang (2025), the ordering appears under FGSM and disappears under AutoAttack, which drops every non-robust model to near zero; on frozen CLIP encoders FGSM reports $16$--$61\%$ robust accuracy that AutoAttack erases.

\textbf{Saturation and the per-image test.} On sixteen frozen CLIP encoders shift-consistency sits at $0.963$--$0.988$ regardless of robustness, so it cannot rank them. Image by image inside one robust encoder, \etaL{} orders each image's radius (Spearman $+0.74$ to $+0.96$) and per-image consistency does not.

\textbf{Ranking robust models by gradient shape.} Among already-robust models \etaL{} no longer separates. The anisotropy $A=\lVert\nabla M\rVert_1/\lVert\nabla M\rVert_2\in[1,\sqrt d]$ does: over thirty RobustBench CIFAR-10 models, lower $A$ is more robust (Spearman $-0.79$; $-0.77$ controlling for clean accuracy), replicated on an MNIST/Fashion training-strength sweep ($-0.61$ to $-0.89$).

\section*{4\quad Mechanism diagnostics}
The margin/Lipschitz decomposition separates two mechanisms: anti-aliasing leaves the margin essentially unchanged and lowers the Lipschitz constant ($\Delta\log L\approx-0.2$), buying robustness through reduced sensitivity; exact cyclic invariance shrinks the margin ($\Delta\log\eta\approx-0.3$ to $-0.4$) and loses robustness, the real-data Ge collapse. For the anisotropy law we propose a spread-over-dimensions deficit $r\approx(\eta/L_1)/(1+CA)$, $C\approx2.65$, verified on a controlled model and thirty real ones. It is a proposed mechanism, not a network-level theorem.

\section*{5\quad What is already done}
\begin{itemize}[leftmargin=1.05em,itemsep=1pt,topsep=2pt]
\item \textbf{Structural theory}, audited and re-verified: the invariant margin identity (Prop.~1), the two-sided bracket (Prop.~2) with a measured $\kappa$, the power-spectrum and bispectrum span lemmas, and the orbit-rank and orbit-flip bounds.
\item \textbf{Capacity-matched dissection} on MNIST, Fashion-MNIST and CIFAR-10, robust radius validated against AutoAttack (mean gap $0.014$), ruling out gradient masking.
\item \textbf{Adversarial-training arm} (PGD-AT, AutoAttack) with the threat-matched \etaL{} law and the shift-consistency anti-prediction.
\item \textbf{Foundation-model regime}: sixteen frozen CLIP encoders, the saturation finding, and the per-image dissociation where training cannot confound.
\item \textbf{Anisotropy study} on thirty RobustBench models with a clean-accuracy control, plus the effective-dimension deficit.
\end{itemize}

\section*{6\quad Limitations and future work}
\begin{itemize}[leftmargin=1.05em,itemsep=1pt,topsep=2pt]
\item \textbf{Not a training target.} \etaL{} is scale-free, so its direct maximizer is the constant classifier: a diagnostic, not a loss.
\item \textbf{The shape law needs range.} On low-dimensional data anisotropy is compressed and does not rank until a training-strength sweep widens the range.
\item \textbf{The deficit is proposed, not proven}, and the cluster-geometry selection theorem (lazy regime) is still open.
\item \textbf{Scale and norms.} Beyond $\ell_\infty$; larger backbones; a genuine text-adversarial axis (current CLIP paraphrase worst case is too mild).
\end{itemize}

\section*{7\quad Project context}
This grew out of a TU Munich practical course that reproduced and extended Ge et al.\ (2021). The supervisor's framing asked for a comparison against Ge et al.\ (2021), the two Kamath et al.\ trade-off papers, and Galloway et al.\ on batch normalization; the dissection reproduces the Galloway effect under fixed-$\epsilon$ accuracy and shows it is threshold-dependent under the radius metric. Repository: \texttt{adversarial-robustness-shift-invariance/}. Compute to date is a modest number of GPU-hours on two RTX A6000.

\vspace{2pt}
\textbf{Selected references.}
{\footnotesize Ge, Singla, Basri, Jacobs, Shift Invariance Can Reduce Adversarial Robustness, NeurIPS 2021.
Kamath, Deshpande, Subrahmanyam, Balasubramanian, Can We Have It All?, NeurIPS 2021.
Tram\`er, Behrmann, Carlini, Papernot, Jacobsen, Fundamental Tradeoffs Between Invariance and Sensitivity, ICML 2020.
Saha, Gokhale, On shift-invariance and adversarial robustness (TIPS), WACV 2024.
Wang et al., Equivariance and adversarial robustness, 2025.
Galloway et al., Batch Normalization Is a Cause of Adversarial Vulnerability, 2019.
Zhang, Making Convolutional Networks Shift-Invariant Again, ICML 2019.
Hein, Andriushchenko, Formal Guarantees on the Robustness of a Classifier, NeurIPS 2017.
Tsuzuku, Sato, Sugiyama, Lipschitz-Margin Training, NeurIPS 2018.
Croce, Hein, Reliable Evaluation with AutoAttack, ICML 2020.
Rony et al., Decoupling Direction and Norm for Min-Norm Attacks, CVPR 2019.
Weng et al., Evaluating the Robustness of Neural Networks: CLEVER, ICLR 2018.
}

\end{multicols}
\end{document}
